\section{Аксиоматика вещественного векторного пространства}

\begin{definition}[Векторное пространство]
Вещественным \textbf{векторным (линейным) пространством} называется непустое множество $V$, элементы которого называются
\emph{векторами}, на котором заданы две операции:
\begin{enumerate}
    \item \textbf{Сложение векторов:} $+ \colon V \times V \to V$, сопоставляющая каждой паре векторов $u, v \in V$
    вектор $u + v \in V$.
    \item \textbf{Умножение вектора на скаляр:} $\cdot \colon \mathbb{R} \times V \to V$, сопоставляющая скаляру
    $\alpha \in \mathbb{R}$ и вектору $v \in V$ вектор $\alpha v \in V$.
\end{enumerate}
При этом для любых $u, v, w \in V$ и любых $\alpha, \beta \in \mathbb{R}$ выполняются следующие 8 аксиом:
\begin{enumerate}
    \item \textbf{Коммутативность сложения:} $u + v = v + u$;
    \item \textbf{Ассоциативность сложения:} $(u + v) + w = u + (v + w)$;
    \item \textbf{Существование нулевого элемента:} существует элемент $0 \in V$ (называемый \emph{нулевым вектором}),
    что $v + 0 = v$ для любого $v \in V$;
    \item \textbf{Существование противоположного элемента:} для каждого $v \in V$ существует элемент $-v \in V$
    такой, что $v + (-v) = 0$;
    \item \textbf{Унитарность умножения на скаляр:} $1 \cdot v = v$;
    \item \textbf{Ассоциативность умножения на скаляры:} $\alpha (\beta v) = (\alpha \beta) v$;
    \item \textbf{Дистрибутивность относительно сложения векторов:} $\alpha (u + v) = \alpha u + \alpha v$;
    \item \textbf{Дистрибутивность относительно сложения скаляров:} $(\alpha + \beta) v = \alpha v + \beta v$.
\end{enumerate}
\end{definition}

\begin{proposition}[Базовые следствия аксиом]
Из приведённых аксиом строго следуют свойства:
\begin{enumerate}
    \item Нулевой вектор $0 \in V$ единственен.
    \item Для каждого вектора $v$ противоположный вектор $-v$ единственен.
    \item Для любого $v \in V$ выполнено $0 \cdot v = 0_V$.
    \item Для любого $\alpha \in \mathbb{R}$ выполнено $\alpha \cdot 0_V = 0_V$.
    \item $(-1) \cdot v = -v$.
    \item Если $\alpha v = 0$, то либо $\alpha = 0$, либо $v = 0$.
\end{enumerate}
\end{proposition}

\begin{example}[Примеры векторных пространств]
    \begin{enumerate}
        \item Пространство $\mathbb{R}^n$ упорядоченных вещественных $n$-ок чисел со стандартными операциями покомпонентного
        сложения и умножения на скаляр.
        \item Пространство многочленов $\mathbb{P}_n$ степени не выше $n$ от переменной $x$ с вещественными коэффициентами.
        \item Пространство вещественных матриц $\mathbb{R}^{m \times n}$ фиксированного размера с операциями матричного
        сложения и умножения на число.
        \item Пространство непрерывных функций $C[a, b]$ на отрезке.
    \end{enumerate}
\end{example}

\section{Подпространства и линейные комбинации}

\begin{definition}[Subset]
Подмножество $U \subseteq V$ называется \textbf{векторным (линейным) подпространством} пространства $V$ (обозначается $U
\le V$), если $U$ само является векторным пространством относительно операций, индуцированных из $V$.
\end{definition}

\begin{theorem}[Критерий подпространства]
Непустое подмножество $U \subseteq V$ является подпространством тогда и только тогда, когда выполнены следующие условия:
\begin{enumerate}
    \item $0 \in U$ (содержит нулевой вектор);
    \item $\forall u, v \in U \implies u + v \in U$ (\textbf{замкнутость относительно сложения});
    \item $\forall u \in U, \; \forall \alpha \in \mathbb{R} \implies \alpha u \in U$ (\textbf{замкнутость относительно
    умножения на скаляр}).
\end{enumerate}
Эквивалентно: $\forall u, v \in U$ и $\forall \alpha, \beta \in \mathbb{R}$ выполнено $\alpha u + \beta v \in U$.
\end{theorem}

\begin{remark}[Важное отличие: линейные vs аффинные подпространства]
Плоскость или прямая в $\mathbb{R}^3$, не проходящая через начало координат (например, $\{ (x_1, x_2, x_3) \mid x_1 +
x_2 + x_3 = 1 \}$), \textbf{не является} векторным подпространством, поскольку не содержит нулевой вектор $0 = (0,0,0)$
и не замкнута относительно умножения на скаляр ($0 \cdot x = 0 \notin U$).
Такие множества называются \emph{аффинными подпространствами} (сдвинутыми подпространствами).
\end{remark}

\begin{definition}[Линейная комбинация и линейная оболочка]
Пусть $S = \{v_1, v_2, \dots, v_k\} \subset V$.
Вектор вида \[ v = c_1 v_1 + c_2 v_2 + \dots + c_k v_k = \sum_{i=1}^k c_i v_i, \quad c_i \in \mathbb{R}, \]
называется \textbf{линейной комбинацией} векторов множества $S$.

Множество всех всевозможных линейных комбинаций векторов из $S$ называется \textbf{линейной оболочкой} (\textbf{span})
множества $S$ и обозначается:
\[ \spanop(S) = \spanop(v_1, \dots, v_k) = \left\{ \sum_{i=1}^k c_i v_i \;\middle|\; c_i \in \mathbb{R} \right\}. \]
Линейная оболочка любого подмножества $S \subseteq V$ всегда является подпространством в $V$.
Множество $S$ называется \textbf{порождающей системой} (generating set) пространства $U$, если $\spanop(S) = U$.
\end{definition}

\section{Линейная независимость, базис и размерность}

\begin{definition}[Линейная независимость]
Система векторов $\{v_1, v_2, \dots, v_k\} \subset V$ называется \textbf{линейно независимой}, если равенство
\[ c_1 v_1 + c_2 v_2 + \dots + c_k v_k = 0 \]
выполняется тогда и только тогда, когда все коэффициенты тривиальны:
\[ c_1 = c_2 = \dots = c_k = 0. \]
Если же существуют такие коэффициенты $c_1, \dots, c_k$, среди которых хотя бы один отличен от нуля ($ \exists i \colon
c_i \ne 0 $), при которых линейная комбинация равна нулю, система векторов называется \textbf{линейно зависимой}.
\end{definition}

\begin{proposition}[Свойство линейной зависимости]
Система векторов $\{v_1, \dots, v_k\}$ ($k \ge 2$) линейно зависима тогда и только тогда, когда хотя бы один из векторов
системы является линейной комбинацией остальных векторов.
\end{proposition}

\begin{definition}[Базис и размерность]
Упорядоченный набор векторов $\mathcal{B} = (e_1, e_2, \dots, e_n)$ пространства $V$ называется \textbf{базисом}
пространства $V$, если:
\begin{enumerate}
    \item Векторы $e_1, \dots, e_n$ линейно независимы;
    \item Система векторов $e_1, \dots, e_n$ порождает всё пространство $V$, то есть $\spanop(e_1, \dots, e_n) = V$.
\end{enumerate}
Число векторов в базисе называется \textbf{размерностью} пространства $V$ и обозначается $\dim V = n$.
Если базис конечен, пространство называется \emph{конечномерным}.
\end{definition}

\begin{theorem}[О единственности координат в базисе]
Пусть $\mathcal{B} = (e_1, \dots, e_n)$ — базис пространства $V$.
Тогда любой вектор $x \in V$ единственным образом представляется в виде линейной комбинации векторов базиса:
\[ x = x_1 e_1 + x_2 e_2 + \dots + x_n e_n. \]
Числа $x_1, \dots, x_n \in \mathbb{R}$ называются \textbf{координатами} вектора $x$ в базисе $\mathcal{B}$, а
координатный столбец обозначается
$[x]_\mathcal{B} = \begin{pmatrix}
                       x_1 \\ \vdots \\ x_n
\end{pmatrix} \in \mathbb{R}^n$.
\end{theorem}
\begin{proof}
    Существование разложения гарантируется тем, что $\mathcal{B}$ порождает $V$ ($\spanop(\mathcal{B}) = V$).
    Докажем единственность от противного.
    Предположим, что вектор $x$ имеет два различных координатных представления в базисе $\mathcal{B}$:
    \[ x = \sum_{i=1}^n x_i e_i \quad \text{и} \quad x = \sum_{i=1}^n y_i e_i. \]
    Вычтем одно равенство из другого:
    \[ 0 = x - x = \sum_{i=1}^n x_i e_i - \sum_{i=1}^n y_i e_i = \sum_{i=1}^n (x_i - y_i) e_i. \]
    Поскольку по определению базиса векторы $e_1, \dots, e_n$ линейно независимы, линейная комбинация может
    равняться нулю только при равенстве нулю всех ее коэффициентов:
    \[ x_i - y_i = 0 \iff x_i = y_i \quad \forall i \in \{1, 2, \dots, n\}. \]
    Таким образом, наборы коэффициентов тождественно совпадают, что доказывает единственность координат.
\end{proof}

\section{Смена базиса и матрица перехода}

Пусть в пространстве $V$ заданы два базиса: «старый» базис $\mathcal{B} = (e_1, \dots, e_n)$ и «новый» базис
$\mathcal{C} = (e'_1, \dots, e'_n)$.
Каждый вектор нового базиса $e'_j$ можно однозначно разложить по векторам старого базиса:
\[ e'_j = \sum_{i=1}^n p_{ij} e_i = p_{1j} e_1 + p_{2j} e_2 + \dots + p_{nj} e_n. \]
\begin{definition}[Матрица перехода]
Матрица $P = P_{\mathcal{B} \to \mathcal{C}} = [P]_{\mathcal{B} \leftarrow \mathcal{C}} \in \mathbb{R}^{n \times n}$,
столбцами которой являются координаты новых базисных векторов в старом базисе, называется \textbf{матрицей перехода} от
базиса $\mathcal{B}$ к базису $\mathcal{C}$:
\[ P = \begin{pmatrix}
           | & | & & | \\
           [e'_1]_\mathcal{B} & [e'_2]_\mathcal{B} & \dots & [e'_n]_\mathcal{B} \\
           | & | & & |
\end{pmatrix}. \]
Матрица перехода всегда невырождена ($\det P \ne 0$), поскольку ее столбцы линейно независимы.
\end{definition}

\begin{theorem}[Преобразование координат при смене базиса]
Пусть вектор $x \in V$ имеет координаты $[x]_\mathcal{B}$ в базисе $\mathcal{B}$ и координаты $[x]_\mathcal{C}$ в базисе
$\mathcal{C}$.
Тогда:
\[ [x]_\mathcal{B} = P_{\mathcal{B} \leftarrow \mathcal{C}} \, [x]_\mathcal{C}, \quad [x]_\mathcal{C} =
\left(P_{\mathcal{B} \leftarrow \mathcal{C}}\right)^{-1} [x]_\mathcal{B} = P_{\mathcal{C} \leftarrow \mathcal{B}} \,
[x]_\mathcal{B}. \]
\end{theorem}

\section{Дуальные пространства, ковекторы и гиперплоскости}

\begin{definition}[Дуальное пространство и ковектор]
Пусть $V$ — векторное пространство над $\mathbb{R}$.
Линейное отображение $f \colon V \to \mathbb{R}$ называется \textbf{линейным функционалом} (или \textbf{линейной
формой}, или \textbf{ковектором}):
\[ f(\alpha u + \beta v) = \alpha f(u) + \beta f(v) \quad \forall u, v \in V, \; \forall \alpha, \beta \in \mathbb{R}.\]
Множество всех линейных функционалов на $V$ само образует векторное пространство относительно поточечных операций
сложения и умножения на скаляр.
Это пространство называется \textbf{дуальным (сопряжённым)} пространством к $V$ и обозначается $V^* = \mathcal{L}(V,
\mathbb{R})$.
Для конечномерных пространств $\dim V^* = \dim V$.
\end{definition}

\begin{definition}[Дуальный базис]
Пусть $\mathcal{B} = (e_1, \dots, e_n)$ — базис пространства $V$.
Базис $\mathcal{B}^* = (\varepsilon^1, \dots, \varepsilon^n)$ в $V^*$, определяемый условием биортогональности (символ
Кронекера):
\[ \varepsilon^i(e_j) = \delta^i_j = \begin{cases}
                                         1, & i = j, \\ 0, & i \ne j,
\end{cases} \] называется \textbf{дуальным (сопряжённым) базисом}.
\end{definition}

\begin{remark}[Ковекторы как линейные измерительные приборы]
Ковектор $\varepsilon^i$ действует как «измерительный прибор», извлекающий $i$-ю координату вектора.
Если $x = \sum_{j=1}^n x_j e_j$, то:
\[ \varepsilon^i(x) = \varepsilon^i\left(\sum_{j=1}^n x_j e_j\right) = \sum_{j=1}^n x_j \varepsilon^i(e_j) =
\sum_{j=1}^n x_j \delta^i_j = x_i.
\] Следовательно, разложение любого вектора по базису может быть записано через дуальные измерения:
\[ x = \sum_{i=1}^n \varepsilon^i(x) e_i. \]
В матричном представлении вектор $x$ изображается столбцом, а ковектор $f$ — строкой: $f(x) = f x = \sum_i f_i x_i$.
\end{remark}

\begin{proposition}[Геометрия ковектора: нарезка на гиперплоскости]
Ненулевой ковектор $f \in V^*$ нарезает пространство $V$ на семейство параллельных аффинных гиперплоскостей (уровневых
множеств, level sets):
\[ H_c = \{ x \in V \mid f(x) = c \}, \quad c \in \mathbb{R}. \]
При этом:
\begin{enumerate}
    \item При $c = 0$ получаем ядро функционала $\ker f = \{ x \in V \mid f(x) = 0 \}$, которое является векторным
    подпространством коразмерности 1 ($\dim \ker f = n - 1$).
    \item При $c \ne 0$ гиперплоскости $H_c$ параллельны подпространству $\ker f$ и сдвинуты относительно начала
    координат.
\end{enumerate}
\end{proposition}

\section{Иерархия геометрических структур}
    В современной математике евклидова геометрия возникает как результат последовательного обогащения абстрактного
    множества дополнительными математическими структурами:
    \begin{gather*}
        \text{\textbf{Метрическое пространство}} \;\longrightarrow\; \text{\textbf{Нормированное пространство}}
    \;\longrightarrow\;\\
        \text{\textbf{Предгильбертово пространство}} \;\longrightarrow\; \text{\textbf{Евклидово
    пространство}}\\
    \end{gather*}

    \subsection{Метрическое пространство}
    Множество $M$ с функцией расстояния (метрикой) $d \colon M \times M \to \mathbb{R}_{\ge 0}$, удовлетворяющей трём
    аксиомам:
    \begin{enumerate}
        \item Тождество неразличимых: $d(x, y) = 0 \iff x = y$;
        \item Симметрия: $d(x, y) = d(y, x)$;
        \item Неравенство треугольника: $d(x, z) \le d(x, y) + d(y, z)$.
    \end{enumerate}
    В метрическом пространстве нет понятий «направление», «ноль» или «сложение» — есть только расстояния между точками.

    \subsection{Нормированное векторное пространство}
    Векторное пространство $V$, на котором задана \textbf{норма} $\|\cdot\| \colon V \to \mathbb{R}_{\ge 0}$,
    удовлетворяющая аксиомам:
    \begin{enumerate}
        \item Положительная определенность: $\|x\| \ge 0$, причём $\|x\| = 0 \iff x = 0$;
        \item Абсолютная однородность: $\|\alpha x\| = |\alpha| \, \|x\|$ для всех $\alpha \in \mathbb{R}, x \in V$;
        \item Неравенство треугольника (субаддитивность): $\|x + y\| \le \|x\| + \|y\|$.
    \end{enumerate}
    \textbf{Индуцирование метрики:} Всякая норма естественным образом порождает метрику расстояния:
    \[ d(x, y) \coloneqq \|x - y\|. \]
    Обратное неверно: не всякая метрика порождается нормой (например, дискретная метрика $d(x, y) = \begin{cases}
                                                                                                        0, & x=y
                                                                                                        \\ 1, & x \ne y
    \end{cases}$ не удовлетворяет условию однородности).

    \subsection{Пространство со скалярным произведением (предгильбертово)}
    Векторное пространство $V$, оснащенное билинейной операцией $\langle \cdot, \cdot \rangle \colon V \times V \to
    \mathbb{R}$, удовлетворяющей аксиомам:
    \begin{enumerate}
        \item Симметричность: $\langle x, y \rangle = \langle y, x \rangle$;
        \item Линейность по первому аргументу: $\langle \alpha x + \beta y, z \rangle = \alpha \langle x, z \rangle +
        \beta \langle y, z \rangle$;
        \item Положительная определенность: $\langle x, x \rangle \ge 0$, причём $\langle x, x \rangle = 0 \iff x = 0$.
    \end{enumerate}

    \begin{theorem}[Неравенство Коши-Буняковского-Шварца]
Для любых элементов $x, y \in V$ со скалярным произведением справедливо:
\[ |\langle x, y \rangle| \le \sqrt{\langle x, x \rangle} \, \sqrt{\langle y, y \rangle}. \]
    \end{theorem}

    \textbf{Индуцирование нормы:} Всякое скалярное произведение порождает евклидову норму:
    \[ \|x\| \coloneqq \sqrt{\langle x, x \rangle}. \]
    Проверка неравенства треугольника для этой нормы опирается непосредственно на неравенство Коши-Буняковского-Шварца:
    \[ \|x + y\|^2 = \langle x+y, x+y \rangle = \|x\|^2 + 2\langle x, y \rangle + \|y\|^2 \le \|x\|^2 + 2\|x\|\|y\| +
    \|y\|^2 = (\|x\| + \|y\|)^2. \]

    \begin{theorem}[Критерий согласованности нормы со скалярным произведением: тождество параллелограмма]
Норма $\|\cdot\|$ на векторном пространстве $V$ порождается некоторым скалярным произведением тогда и только тогда,
когда для любых $x, y \in V$ выполняется \textbf{тождество параллелограмма}:
\[ \|x + y\|^2 + \|x - y\|^2 = 2\|x\|^2 + 2\|y\|^2. \]
Если тождество параллелограмма выполнено, скалярное произведение однозначно восстанавливается по поляризационному
тождеству:
\[ \langle x, y \rangle = \frac{1}{4}\left(\|x + y\|^2 - \|x - y\|^2\right). \]
    \end{theorem}

    \begin{example}[Контрпример: нормы, не порождаемые скалярным произведением]
Рассмотрим пространство $\mathbb{R}^2$:
\begin{itemize}
    \item Манхэттенская норма ($L_1$): $\|x\|_1 = |x_1| + |x_2|$.
    Возьмём $x = (1, 0)$ и $y = (0, 1)$.
    Тогда $\|x\|_1 = 1$, $\|y\|_1 = 1$.
    Имеем $x+y = (1, 1)$, $x-y = (1, -1)$, отсюда $\|x+y\|_1^2 + \|x-y\|_1^2 = 2^2 + 2^2 = 8$.
    Однако $2\|x\|_1^2 + 2\|y\|_1^2 = 2(1) + 2(1) = 4$.
    Так как $8 \ne 4$, норма $L_1$ не порождается скалярным произведением.
    \item Чебышёвская норма ($L_\infty$): $\|x\|_\infty = \max(|x_1|, |x_2|)$.
    Аналогично не удовлетворяет тождеству параллелограмма.
\end{itemize}
    \end{example}

    \subsection{Евклидово пространство}
    \textbf{Евклидовым пространством} называется конечномерное вещественное векторное пространство, оснащенное скалярным
    произведением.
    В нём одновременно согласованы все уровни геометрии:
    \[ \begin{array}{lcccc}
           \text{Скалярное произведение:} & \langle x, y \rangle & & & \\
           & \Downarrow & & & \\
           \text{Длины (норма):} & \|x\| = \sqrt{\langle x, x \rangle} & & & \\
           & \Downarrow & \searrow & & \\
           \text{Расстояния (метрика):} & d(x, y) = \|x - y\| & & \text{Углы:} & \cos \theta = \frac{\langle x, y
           \rangle}{\|x\| \|y\|}
    \end{array} \]

\section{Ортогональность и ортогональное проецирование}

\begin{definition}[Orthogonality]
Два вектора $u, v \in V$ называются \textbf{ортогональными} ($u \perp v$), если их скалярное произведение равно нулю:
$\langle u, v \rangle = 0$.
\end{definition}

\subsection{Вывод формулы одномерной проекции из ортогональности}
Пусть даны ненулевой вектор $u \in V$ ($u \ne 0$) и произвольный вектор $v \in V$.
Мы ищем ортогональную проекцию вектора $v$ на прямую, натянутую на $u$, то есть вектор $p = P_u(v) \in \spanop(u)$.
\begin{enumerate}
    \item Так как $p \in \spanop(u)$, существует скаляр $\alpha \in \mathbb{R}$ такой, что $p = \alpha u$.
    \item Определим вектор ошибки (перпендикуляр) $e \coloneqq v - p = v - \alpha u$.
    \item Требование \textbf{ортогональности проекции} состоит в том, чтобы вектор ошибки был перпендикулярен
    направляющему вектору $u$:
    \[ e \perp u \iff \langle e, u \rangle = 0 \iff \langle v - \alpha u, u \rangle = 0. \]
    \item Раскрывая скалярное произведение по линейности:
    \[ \langle v, u \rangle - \alpha \langle u, u \rangle = 0 \iff \alpha \langle u, u \rangle = \langle v, u
    \rangle. \]
    \item Поскольку $u \ne 0$, имеем $\langle u, u \rangle = \|u\|^2 > 0$.
    Отсюда однозначно находим коэффициент $\alpha$:
    \[ \alpha = \frac{\langle v, u \rangle}{\langle u, u \rangle} = \frac{\langle v, u \rangle}{\|u\|^2}. \]
    \item Подставляя $\alpha$ в выражение для $p$, получаем каноническую \textbf{формулу ортогональной проекции}:
    \[ P_u(v) = \frac{\langle v, u \rangle}{\langle u, u \rangle} u = \frac{\langle v, u \rangle}{\|u\|^2} u. \]
\end{enumerate}

\subsection{Проекция на многомерное подпространство}
Пусть $U \le V$ — подпространство с базисом, составленным из столбцов матрицы $A\in \mathbb{R}^{n \times k}$ полного
столбцового ранга ($\rank A = k$).
Тогда проекция вектора $y \in \mathbb{R}^n$ на $U = \Col(A)$ представляется в виде $p = A x$ для некоторого $x \in
\mathbb{R}^k$.
Условие ортогональности ошибки $e = y - Ax$ всем столбцам матрицы $A$ записывается в виде:
\[ A^\top (y - Ax) = 0 \iff A^\top A x = A^\top y. \]
Так как столбцы $A$ линейно независимы, матрица Грама $A^\top A \in \mathbb{R}^{k \times k}$ симметрична и строго
положительно определёна, а значит, обратима.
Находим $x$:
\[ x = (A^\top A)^{-1} A^\top y. \]
Тогда проекция равна $p = A x = A (A^\top A)^{-1} A^\top y$.
Отсюда получаем \textbf{матрицу ортогонального проектора}:
\[ P = A (A^\top A)^{-1} A^\top. \]
Матрица $P$ обладает фундаментальными свойствами:
\begin{enumerate}
    \item \textbf{Идемпотентность:} $P^2 = P$ (повторное проецирование ничего не меняет);
    \item \textbf{Симметричность:} $P^\top = P$ (гарантирует ортогональность проецирования).
\end{enumerate}

\section{Метод главных компонент (PCA)}

\textbf{Метод главных компонент (Principal Component Analysis, PCA)} — один из ключевых методов снижения размерности
данных без учителя в машинном обучении и математической статистике.

\subsection{Постановка задачи и алгоритм}
Пусть дана матрица центрированных наблюдений $X \in \mathbb{R}^{n \times d}$, где $n$ — число объектов, $d$ — число
исходных признаков, причём выборочное среднее каждого признака равно нулю ($\sum_{i=1}^n x_{ij} = 0$).
Выборочная ковариационная матрица признаков задаётся формулой:
\[ \Sigma = \frac{1}{n - 1} X^\top X \in \mathbb{R}^{d \times d}. \]
Матрица $\Sigma$ симметрична и положительно полуопределена.

Цель PCA — найти такое ортогональное направление $w_1 \in \mathbb{R}^d$ ($\|w_1\|_2 =1$), проекция данных на которое
обладает максимальной выборочной дисперсией:
\[ \max_{\|w_1\| = 1} \operatorname{Var}(X w_1) = \max_{\|w_1\| = 1} \frac{1}{n-1} (X w_1)^\top (X w_1) =
\max_{\|w_1\| = 1} w_1^\top \Sigma w_1. \]
Используя метод множителей Лагранжа для функции $L(w_1, \lambda) =w_1^\top \Sigma w_1 - \lambda (w_1^\top w_1 - 1)$,
получаем необходимое условие экстремума:
\[ \nabla_{w_1} L = 2 \Sigma w_1 - 2 \lambda w_1 = 0 \iff \Sigma w_1 = \lambda w_1. \]
Таким образом, направление первой главной компоненты $w_1$ — это \textbf{собственный вектор} ковариационной матрицы
$\Sigma$, отвечающий ее максимальному собственному значению $\lambda_1$.
Дисперсия вдоль этого направления равна в точности $\lambda_1$.

Последующие главные компоненты $w_2,\dots,w_k$ находятся как собственные векторы $\Sigma$, ортогональные предыдущим,
упорядоченные по убыванию собственных значений $\lambda_1 \ge \lambda_2 \ge \dots \ge \lambda_d \ge 0$.

\subsection{Связь PCA с сингулярным разложением (SVD)}
На практике вычисление ковариационной матрицы $\Sigma = \frac{1}{n-1} X^\top X$ не требуется.
Достаточно применить SVD непосредственно к центрированной матрице данных:
\[ X = U S V^\top. \]
Тогда:
\[ X^\top X = (V S U^\top)(U S V^\top) = V S^2 V^\top. \]
Отсюда видно, что:
\begin{enumerate}
    \item Правые сингулярные векторы $V = (v_1, \dots, v_d)$ являются в точности главными компонентами (собственными
    векторами матрицы ковариации $\Sigma$).
    \item Собственные значения ковариационной матрицы связаны с сингулярными числами матрицы $X$ соотношением:
    \[ \lambda_i = \frac{s_i^2}{n - 1}. \]
    \item Новые низкоразмерные координаты данных (проекции) вычисляются умножением: $Z_k = X V_k = U_k S_k$.
\end{enumerate}

\newpage

% =========================================================================
% ГЛАВА 2
% =========================================================================